\ifx\allfiles\undefined
%生成目录超链接，使用UTF8字符集
\documentclass[12pt,UTF8]{ctexart}


\usepackage{mathtools}%花括号括起的公式
\usepackage{enumerate}%列表
\usepackage{graphicx}%表格
\usepackage[table]{xcolor}%彩色表格
\usepackage{booktabs}%调整表格行高


\renewcommand\tabcolsep{50pt}
\renewcommand\arraystretch{2}

%文本框
\usepackage{tcolorbox}
\tcbuselibrary{skins,vignette,raster,listings,listingsutf8,theorems, breakable,magazine,fitting}

%页面设置
\usepackage{geometry}%设置页面
\usepackage{setspace}%设置局部行距

\usepackage{ctex}
\usepackage{CJK}

%图像处理
\usepackage{graphicx}
\usepackage{float}%允许浮动
\usepackage{caption}%设置标题
\usepackage{graphicx, subfig}

%附录
\usepackage[toc,page]{appendix}

%导言区
\title{中小微企业的信贷决策}

%不要页眉
\pagestyle{plain}
\bibliographystyle{plain}

\geometry{a4paper}%将纸张设置为A4大小

%修改文档结构层次的名称
\usepackage{titlesec}%chapter1修改为一
\titleformat{\part}{\centering\Huge\bfseries}{\thepart}{1em}{}
\renewcommand\thepart{\chinese{part}}

%subsubsubsection的使用
%\usepackage{titlesec}
%\titleclass{\subsubsubsection}{straight}[\subsection]
%
%\newcounter{subsubsubsection}[subsubsection]
%\renewcommand\thesubsubsubsection{\thesubsubsection.\arabic{subsubsubsection}}
%\renewcommand\theparagraph{\thesubsubsubsection.\arabic{paragraph}}
%
%\titleformat{\subsubsubsection}
%{\normalfont\normalsize\bfseries}{\thesubsubsubsection}{1em}{}
%\titlespacing*{\subsubsubsection}
%{0pt}{3.25ex plus 1ex minus .2ex}{1.5ex plus .2ex}
%
%\makeatletter
%\renewcommand\paragraph{\@startsection{paragraph}{5}{\z@}%
%	{3.25ex \@plus1ex \@minus.2ex}%
%	{-1em}%
%	{\normalfont\normalsize\bfseries}}
%\renewcommand\subparagraph{\@startsection{subparagraph}{6}{\parindent}%
%	{3.25ex \@plus1ex \@minus .2ex}%
%	{-1em}%
%	{\normalfont\normalsize\bfseries}}
%\def\toclevel@subsubsubsection{4}
%\def\toclevel@paragraph{5}
%\def\toclevel@paragraph{6}
%\def\l@subsubsubsection{\@dottedtocline{4}{7em}{4em}}
%\def\l@paragraph{\@dottedtocline{5}{10em}{5em}}
%\def\l@subparagraph{\@dottedtocline{6}{14em}{6em}}
%\makeatother
%
%\setcounter{secnumdepth}{4}

\begin{document}
\else
\fi
%分文件开始

\appendix
\renewcommand{\appendixpagename}{附录}
\renewcommand{\appendixtocname}{附录}
%附录开始
\begin{appendices}

\noindent \textbf{{\large 支撑材料清单：}}
\begin{enumerate}[(1)]
	\item 问题1多元线性规划求解代码：Q1LP.mlx 
	\item 问题2多元线性规划求解源代码：Q2LP.mlx 
	\item 问题3多元线性规划求解源代码：Q3LP.mlx 
	\item 问题1、2层次分析法源代码：Q2Work.mlx 
	\item 问题3层次分析法源代码：Q3Work.mlx
	\item 多元线性回归源数据：附件1分档.xIsx
	\item 多元线性回归源数据：附件2分档.xIsx
	\item 多元线性回归源数据：问题3附件2分档.xlsx
	\item 规划求解结果.xlsx
	\item 决策树训练结果.pdf
	\item PowerBI数据处理.pbix
	\item 问题1数据处理.xIsx
	\item 问题2数据处理.xIsx
	\item 问题3数据处理.xIsx
\end{enumerate}

	\section{问题一中层次分析法的MATLAB实现}

	首先编写实现单层次结构的层次分析法的\verb|AHP(A)|函数
	\begin{tcolorbox}[breakable]
		\begin{verbatim}
		function [W, Lmax, CI, CR] = AHP(A)
		% 实现单层次结构的层次分析法
		% 输入: A为成对比较矩阵
		% 输出: W为权重向量, Lmax为最大特征值, CI为一致性指标, CR为一致性比率
		
		[V,D] = eig(A);
		[Lmax,ind] = max(diag(D));       % 求最大特征值及其位置
		W = V(:,ind) / sum(V(:,ind));    % 最大特征值对应的特征向量做标准化
		Lmax = mean((A * W) ./ W);       % 计算最大特征值
		n = size(A, 1);                  % 矩阵行数
		CI = (Lmax - n) / (n - 1);       % 计算一致性指标
		% Saaty随机一致性指标值
		RI = [0 0 0.58 0.90 1.12 1.24 1.32 1.41 1.45 1.49 1.51]; 
		CR = CI / RI(n);                 % 计算一致性比率
	 \end{verbatim}
	\end{tcolorbox}

	按照如下代码构建目标层与准则层之间的成对比较矩阵，并调用\verb|AHP(A)|，即可得到组合权向量：

	\begin{tcolorbox}[breakable]
		\begin{verbatim}
		%目标层与准则层之间的成对比较矩阵
		A =
		[1  8  3  5  6  1; 
			1/8  1  1/7  1/5  1/4  1/8;
			 1/3  7  1  3  2  1/3;
			 1/5  5  1/3  1  1/3  1/5;
			 1/6  4  1/2  3  1  1/4;
			 	1  8  3  5  4  1;];
		[W, Lmax, CI, CR] = AHP(A)
	\end{verbatim}
	\end{tcolorbox}


	\section{问题一中规划模型的MALTAB实现}
	\begin{tcolorbox}[breakable]
		\begin{verbatim}
%Part1:导入数据
%此处xlsread函数的参数含义:路径，读取第一个工作表，读取单元格范围
%年利率
AnnualInterestRate = xlsread("D:\OneDrive\GitHub\CUMCM2020材料包\附件1分档.xlsx",1,'F11:F133');
%保有率
RetentionRate = xlsread("D:\OneDrive\GitHub\CUMCM2020材料包\附件1分档.xlsx",1,'G11:G133');
%最大贷款额度
MaxLoanLimit = xlsread("D:\OneDrive\GitHub\CUMCM2020材料包\附件1分档.xlsx",1,'H11:H133');



%Part2:运算台
%xi的权重向量ki=该企业的保有率*年利率*1年
k=AnnualInterestRate'.*RetentionRate'.*1;
%企业数
n=123;
%银行信贷总额，单位万元
total=3050:100:5000;
num=size(total,2);

%Part3:规划求解
%目标函数中，xi的权重向量f
%目标函数应标准化求解最小值
f=-k;

%不等式约束条件A,Ax<=b，inequality constraint
A=[eye(n) ;ones(1,n)];

%等式约束条件aeqx=beq，equality constraint
aeq=zeros(1,n);
beq=0;

%xi的上下界，bound constraint
lb=zeros(n,1);
ub=100*ones(n,1);



%Part4:结果台
%预分配内存
income=zeros(1,num);
test=zeros(n,num);

%构建循环
for i = 1:1:num
b=[MaxLoanLimit;total(i)];

%各家企业实际贷款额度
x=linprog(f,A,b,aeq,beq,lb,ub);

%银行收入
%目标函数反向标准化，求解最大值
income(i)=sum(-f.*x,'all');

test(:,i)=x;
end

plot(total,income)
test
	\end{verbatim}
	\end{tcolorbox}



	\section{问题二中决策树模型的工程参数}
	\begin{center}
		\begin{table}[H]  %table 里面也可以嵌套tabular,只有tabular是不能加标题的
			\centering
			\resizebox{300pt}{200pt}{
				\begin{tabular}{cc}
					\hline
					参数名                           & 参数值 \\
					\hline
					数据切分                         & 0.75   \\
					数据洗牌                         & 是     \\
					交叉验证                         & 否     \\
					随机种子                         & 0      \\
					分裂准则                         & mse    \\
					树的最大深度                     & None   \\
					分裂节点时的策略                 & best   \\
					分裂一个内部节点所需最少样本数   & 2      \\
					每个叶子节点所需最少样本数       & 2      \\
					叶子节点中样本的最小权重         & 0      \\
					搜寻最佳划分的时候考虑的特征数量 & None   \\
					叶子节点的最大数量               & None   \\
					信息墒或者基尼系数不纯度的阀值   & 0      \\
					\hline
				\end{tabular}}
		\end{table}
	\end{center}

	\section{问题二中规划模型的MATLAB实现}
	\begin{tcolorbox}[breakable]
		\begin{verbatim}
		%Part1:导入数据
		%此处xlsread函数的参数含义:路径，读取第一个工作表，读取单元格范围
		%年利率
		AnnualInterestRate = xlsread("D:\OneDrive\GitHub\CUMCM2020材料包\附件2分档.xlsx",1,'F11:F312');
		%保有率
		RetentionRate = xlsread("D:\OneDrive\GitHub\CUMCM2020材料包\附件2分档.xlsx",1,'G11:G312');
		%最大贷款额度
		MaxLoanLimit = xlsread("D:\OneDrive\GitHub\CUMCM2020材料包\附件2分档.xlsx",1,'H11:H312');
		
		
		
		%Part2:运算台
		%xi的权重向量ki=该企业的保有率*年利率*1年
		k=AnnualInterestRate'.*RetentionRate'.*1;
		%企业数
		n=302;
		%银行信贷总额，单位万元
		total=10000;
		
		
		
		%Part3:规划求解
		%目标函数中，xi的权重向量f
		%目标函数应标准化求解最小值
		f=-k
		
		%不等式约束条件Ax<=b，inequality constraint
		A=[eye(n) ;ones(1,n)];
		b=[MaxLoanLimit;total];
		
		%等式约束条件aeqx=beq，equality constraint
		aeq=zeros(1,n);
		beq=0;
		
		%xi的上下界，bound constraint
		lb=zeros(n,1);
		ub=100*ones(n,1);
		
		
		
		%Part4:结果台
		%各家企业实际贷款额度
		x=linprog(f,A,b,aeq,beq,lb,ub)
		
		%银行收入，单位万元
		%目标函数反向标准化，求解最大值
		income=sum(-f.*x,'all')
	\end{verbatim}
	\end{tcolorbox}


	\section{问题三中层次分析法的MATLAB实现}
	\begin{tcolorbox}[breakable]
		\begin{verbatim}
A = 
[1  5  7  9  3;
1/5  1  3  1/5  1;
1/7  1/3  1  5  1/5;
1/9  1/5  1/5  1  1/7;
1/3  1  5  7  1;];
[W, Lmax, CI, CR] = AHP(A)
	\end{verbatim}
	\end{tcolorbox}

	\section{问题三中规划模型的MATLAB实现}
	\begin{tcolorbox}[breakable]
		\begin{verbatim}
%Part1:导入数据
%此处xlsread函数的参数含义:路径，读取第一个工作表，读取单元格范围
%年利率
AnnualInterestRate = xlsread("D:\OneDrive\GitHub\CUMCM2020材料包\问题3附件2分档.xlsx",1,'F11:F312');
%保有率
RetentionRate = xlsread("D:\OneDrive\GitHub\CUMCM2020材料包\问题3附件2分档.xlsx",1,'G11:G312');
%最大贷款额度
MaxLoanLimit = xlsread("D:\OneDrive\GitHub\CUMCM2020材料包\问题3附件2分档.xlsx",1,'H11:H312');
%问题2求解结果
Q2Result = xlsread("D:\OneDrive\GitHub\CUMCM2020材料包\规划求解结果.xlsx",1,'S2:S303');



%Part2:运算台
%xi的权重向量ki=该企业的保有率*年利率*1年
k=AnnualInterestRate'.*RetentionRate'.*1;
%企业数
n=302;
%银行信贷总额，单位万元
total=10000;



%Part3:规划求解
%目标函数中，xi的权重向量f
%目标函数应标准化求解最小值
f=-k

%不等式约束条件Ax<=b，inequality constraint
A=[eye(n) ;ones(1,n)];
b=[MaxLoanLimit;total];

%等式约束条件aeqx=beq，equality constraint
aeq=zeros(1,n);
beq=0;

%xi的上下界，bound constraint
lb=zeros(n,1);
ub=100*ones(n,1);



%Part4:结果台
%各家企业实际贷款额度
x=linprog(f,A,b,aeq,beq,lb,ub)

%银行收入
%目标函数反向标准化，求解最大值
income=sum(-f.*x,'all')

Delta=x-Q2Result
	\end{verbatim}
	\end{tcolorbox}


	%附录结束
\end{appendices}


%分文件结束
\ifx\allfiles\undefined
\end{document}
\fi